\documentclass[aspectratio=169, 10pt, handout]{beamer}

%%%%%%%%%%%%%%%%%%%
% Beamer settings %
%%%%%%%%%%%%%%%%%%%
\usefonttheme{serif}
\setbeamertemplate{footline}[frame number]{}
\setbeamertemplate{navigation symbols}{}
\setbeamertemplate{items}[circle]
\setbeamertemplate{section in toc}[sections numbered]
\setbeamertemplate{frametitle continuation}{}
\setbeamertemplate{caption}[numbered]
\setbeamertemplate{section in toc shaded}[default][100]
\setbeamertemplate{enumerate items}[default]
\setbeamerfont{section in toc shaded}{series=\mdseries}
\setbeamerfont{section in toc}{series=\bfseries}

%%%%%%%%%%
% Images %
%%%%%%%%%%
\usepackage{graphicx}

%%%%%%%%%%
% Tables %
%%%%%%%%%%
\usepackage{booktabs}
\usepackage{multirow}

%%%%%%%%%%%%%%%%%%
% Math: Pacakges %
%%%%%%%%%%%%%%%%%%
\usepackage{amsmath}
\usepackage{amsfonts}
\usepackage{amssymb}
\usepackage{amsthm}
\usepackage{bm}
\usepackage{mathrsfs}

%%%%%%%%%%%%%%%%%%%%%%%%%%%%%%
% Math: Theorem Environments %
%%%%%%%%%%%%%%%%%%%%%%%%%%%%%%
\setbeamertemplate{theorems}[numbered]

\theoremstyle{definition}
\newtheorem{defn}{Definition}
\newtheorem*{defn*}{Definition}
\newtheorem{asm}{Assumption}
\newtheorem*{asm*}{Assumption}

\theoremstyle{plain}
\newtheorem{thm}{Theorem}
\newtheorem{lem}{Theorem}
\newtheorem{cor}{Corollary}
\newtheorem*{thm*}{Theorem}
\newtheorem*{lem*}{Lemma}
\newtheorem*{cor*}{Corollary}

%%%%%%%%%%%%%%%%%%%%%%%%%%%%%%%
% Math: Operators and Symbols %
%%%%%%%%%%%%%%%%%%%%%%%%%%%%%%%
% Does not imply
% Source - https://tex.stackexchange.com/a/509385
% Posted by jds
% Retrieved 2026-03-30, License - CC BY-SA 4.0
\newcommand{\notimplies}{\;\not\mspace{-12mu}\implies}

% Independence in Probability Theory and Statistics.
\newcommand{\indep}{\mathrel{\perp\mspace{-11mu}\perp}}

% Indicators
\def\ind{\mathbf{1}}

\DeclareMathOperator*{\argmax}{argmax}

% Probability operators
\def\Pr{\mathrm{Pr}}
\def\E{\mathrm{E}}
\def\Var{\mathrm{Var}}
\def\Cov{\mathrm{Cov}}
\def\Supp{\mathrm{support}}
\def\pto{\overset{\mathrm{p}}{\to}}
\def\dto{\rightsquigarrow}
\def\Normal{\mathrm{N}} % Normal distribution
\def\Unif{\mathrm{Unif}} % Normal distribution

\def\Op{O_{\mathrm{p}}}
\def\op{o_{\mathrm{p}}}

% IID
\def\iid{\mathrm{iid}}

\def\sgn{\mathrm{sgn}}

\def\bfQ{\mathbf{Q}}
\def\R{\mathbb{R}}
\def\N{\mathbb{N}}
\def\bW{\mathbf{W}}
\def\bY{\mathbf{Y}}
\def\bA{\mathbf{A}}
\def\bZ{\mathbf{Z}}
\def\bX{\mathbf{X}}
\def\bU{\mathbf{U}}
\def\calA{\mathcal{A}}
\def\calS{\mathcal{S}}
\def\calZ{\mathcal{Z}}
\def\calY{\mathcal{Y}}
\def\calW{\mathcal{W}}
\def\calQ{\mathcal{Q}}
\def\d{\mathrm{d}}
\def\sphere{\mathbb{S}}

\newcommand\mapsfrom{\mathrel{\reflectbox{\ensuremath{\mapsto}}}}

% \usepackage[
% authordate,backend=biber,natbib,maxcitenames=2
% ]{biblatex-chicago}

% \addbibresource{../main.bib}

\makeatletter
\renewcommand<>\appendix{%
\setcounter{section}{0}%
\renewcommand\insertsectionnumber{\Alph{section}}%
  \only#1{\part{\appendixname}
  \addtocontents{nav}{\protect\headcommand{\protect\beamer@appendixpages{\the\c@page}}}}}
\makeatother

\providecommand{\mainframetotal}{?}
\providecommand{\appendixframetotal}{?}
\newif\ifappendixstarted

\setbeamertemplate{footline}{%
  \hfill
  \usebeamerfont{page number in head/foot}%
  \usebeamercolor[fg]{page number in head/foot}%
  \insertframenumber/\mainframetotal\kern1em\vskip2pt%
}

\makeatletter
\AddToHook{cmd/appendix/before}{%
  \immediate\write\@auxout{%
    \string\gdef\string\mainframetotal{\arabic{framenumber}}%
  }%
}

\AddToHook{cmd/appendix/after}{%
  \global\appendixstartedtrue
  \setcounter{framenumber}{0}%
  \setbeamertemplate{footline}{%
    \hfill
    \usebeamerfont{page number in head/foot}%
    \usebeamercolor[fg]{page number in head/foot}%
    A\insertframenumber/A\appendixframetotal\kern1em\vskip2pt%
  }%
}

\AtEndDocument{%
  \ifappendixstarted
    \immediate\write\@auxout{%
      \string\gdef\string\appendixframetotal{\arabic{framenumber}}%
    }%
  \else
    \immediate\write\@auxout{%
      \string\gdef\string\mainframetotal{\arabic{framenumber}}%
    }%
  \fi
}
\makeatother

\title{A probit model with a continuous endogenous regressor}
\author{Ahnaf Rafi}
\date{\today \\
ECON 8710: Cross Section Econometrics \\
University of Virginia, Fall 2026
}
\begin{document}

\begin{frame}[plain,noframenumbering]
\titlepage
\end{frame}

\section*{Introduction}

\begin{frame}{Observables and the structural model}
We observe an i.i.d.\ sample of \(\xi = (Y, D, X', Z')'\).
\(Y \in \{0, 1\}\) is the binary outcome,
\(D \in \R\) is a continuous \textbf{endogenous} covariate,
\(X \in \R^{k}\) are exogenous covariates,
\(Z \in \R^{l}\) are excluded instruments, with \(l \geq 1\).
For unobserved \(\epsilon\) and \(\nu\), model is
\begin{equation*}
  \begin{split}
    Y =
    & \, \ind \{\alpha_{0} D + X' \beta_{0} - \epsilon \geq 0\}, \\
    D =
    & \, X' \gamma_{0} + Z' \pi_{0} + \nu = S' \delta_{0} + \nu, \quad
    S = (X', Z')', \quad \delta_{0} = (\gamma_{0}', \pi_{0}')'.
  \end{split}
\end{equation*}
Assume \(\pi \neq \mathbf{0}_{l}\) (rank condition), \(\E [S S']\) is finite and
positive definite, \((\epsilon, \nu) \indep S\) and
\begin{equation*}
  \begin{pmatrix} \epsilon \\ \nu \end{pmatrix}
  \sim \Normal (\mathbf{0}_{2}, \Sigma_{0}), \quad
  \Sigma_{0} =
  \begin{pmatrix}
    1 & \rho_{0} \sqrt{\tau_{0}} \\
    \rho_{0} \sqrt{\tau_{0}} & \tau_{0}
  \end{pmatrix}.
\end{equation*}
\vfill
\begin{itemize}
  \item We require that \(\rho_{0} \in [-1, 1]\).
    With this parametrization, \(\Sigma_{0}\) is a covariance matrix.
    \vfill
  \item When \(\rho_{0} \neq 0\), \(D\) is endogenous since it correlates with
    \(\varepsilon\).
    \vfill
  \item The full vector of unknown parameters is \(\vartheta_{0} = (\alpha_{0},
    \beta_{0}', \gamma_{0}', \pi_{0}', \tau_{0}, \rho_{0})\).
\end{itemize}
\end{frame}

\begin{frame}{Conditional log-likelihood}
Let \(p_{\vartheta} (\cdot | \cdot)\) be the conditional density of \((Y,
D)\) given \(S\) at candidate \(\vartheta\) for \(\vartheta_{0}\).
For \(\tau > 0\) and \(|\rho| < 1\),
\begin{equation*}
  \begin{split}
    p_{\vartheta} (Y, D | S) =
    & \, \frac{\exp \left\{-\frac{(D - S' \delta)^{2}}{2 \tau}\right\}}
    {\sqrt{2 \pi \tau}} \times \Phi \left(
    \frac{\sqrt{\tau} (\alpha D + X' \beta) - \rho (D - S' \delta)}
    {\sqrt{\tau (1 - \rho^{2})}}\right)^{Y} \\
    & \times \Phi \left(
    \frac{\rho (D - S' \delta) - \sqrt{\tau} (\alpha D + X' \beta)}
    {\sqrt{\tau (1 - \rho^{2})}}\right)^{1 - Y}.
  \end{split}
\end{equation*}
Take logarithms:
\begin{equation*}
  \begin{split}
    \log p_{\vartheta} (Y, D | S) =
    & \, -\frac{\log (2 \pi)}{2} - \frac{1}{2} \log \tau
    - \frac{(D - S' \delta)^{2}}{2 \tau} \\
    & + Y \log \Phi \left(
    \frac{\sqrt{\tau} (\alpha D + X' \beta) - \rho (D - S' \delta)}
    {\sqrt{\tau (1 - \rho^{2})}}\right) \\
    & + (1 - Y) \log \Phi \left(
    \frac{\rho (D - S' \delta) - \sqrt{\tau} (\alpha D + X' \beta)}
    {\sqrt{\tau (1 - \rho^{2})}}\right).
  \end{split}
\end{equation*}
\end{frame}

\begin{frame}{Conditional log-likelihood}
We can drop the constant term and use \(2 Y - 1 \in \{-1, 1\}\) to simplify
things a little:
\begin{equation*}
  \ell_{\vartheta} (\xi) = -\frac{1}{2} \log \tau
  - \frac{(D - S' \delta)^{2}}{2 \tau}
  + \log \Phi \left((2 Y - 1)
  \frac{\sqrt{\tau} (\alpha D + X' \beta) - \rho (D - S' \delta)}
  {\sqrt{\tau (1 - \rho^{2})}}\right).
\end{equation*}
This will also help reduce computational load later on.
The conditional log-likelihood is (up to an additive constant) equal to the
sample sum of the above;
the negative average conditional log-likelihood is that multiplied by a factor
of \(-1/n\).
\end{frame}

\end{document}
